\documentclass[12pt]{article}
\usepackage[utf8]{inputenc}
\usepackage[T1]{fontenc}
\usepackage{hyperref}
\title{American Sign Language Language Resources for K-12 Curriculum Development}
\author{Compiled by Claude Sonnet 5.5 (claude-sonnet-5-5) for Andr\'as Kornai}
\date{}
\begin{document}
\maketitle
\section{Introduction}
American Sign Language (ISO 639-3 ase) is a Deaf community sign language of the United States. The search found mostly teacher-facing curriculum documents; no current official U.S. national K-12 ASL standards document turned up.
\section{Primary Educational Resources}
\begin{itemize}
\item Saskatchewan (Canada) resource list for ASL and Deaf Culture 10, 20, 30: \url{https://curriculum.gov.sk.ca/FullResourceList?id=599}
\item MDElio ASL curriculum, instruction and assessment materials (subscription): \url{https://mdelio.org/node/861}
\item New York State world languages guide for American Sign Language (PDF): \url{https://www.nysed.gov/sites/default/files/programs/world-languages/amsign.pdf}
\item Manitoba K-12 ASL framework (only the acknowledgments page was seen): \url{https://www.edu.gov.mb.ca/k12/cur/languages/asl/framework/acknowledgements.pdf}
\item Education.com ASL worksheets and games (supplementary): \url{https://www.education.com/resources/asl/}
\end{itemize}
\section{Linguistic and Reference Materials}
No dictionaries or grammars were found in these search results.
\section{Religious and Community Resources}
None found in these search results.
\section{Digital and Audio Resources}
See the video and web resources linked from the Saskatchewan list and the subscription site above.
\section{Curriculum Development Context}
Existing materials are provincial or state level and mostly teach ASL as a school subject. A curriculum conducted in ASL for deaf and hearing learners would need to build on these frameworks and verify their current status.
\end{document}